\section*{Chapitre \ref{ch:iv2:rates-and-fx} --- Taux et change}
\markboth{\protect\omTocMark{}{Réponses~: Taux et change}}{\protect\omTocMark{}{Réponses~: Taux et change}}

\begin{solution}{iq:iv2:rates-and-fx:1}
Il vaut environ 82.19 avec une capitalisation annuelle, et environ 81.87 avec une capitalisation continue.

$100/1.04^5 = 82.19$, et $100e^{-0.2} = 81.87$. Un taux affiché ne signifie rien sans sa base de capitalisation.

\iqwrong{Utiliser $100(1 - 5 \times 0.04) = 80$. C'est une règle d'actualisation simple, et elle est fausse de 2 points ici.}
\end{solution}

\begin{solution}{iq:iv2:rates-and-fx:2}
Il vaut environ 100.96. À un rendement de 5~\%, il vaut exactement 100, car une obligation dont le coupon est égal à son rendement se traite au pair.

$105/1.04 = 100.96$. À 5~\%, $105/1.05 = 100$.

\iqwrong{Croire qu'une prime signifie que l'obligation est une mauvaise affaire. Le coupon la paie, et le prix revient au pair à l'échéance.}
\end{solution}

\begin{solution}{iq:iv2:rates-and-fx:3}
La duration de Macaulay vaut 5, et la duration modifiée vaut $5/1.04 \approx 4.81$. Donc 10 points de base font perdre environ 0.481~\%, soit environ 40 cents sur un prix de 82.19.

Un zéro-coupon ne comporte qu'un seul paiement, si bien que sa durée moyenne est son échéance. La variation en pourcentage est la duration modifiée multipliée par la variation du rendement.

\iqwrong{Appliquer les 4.81~\% aux 100~\$ de valeur nominale. Ce taux s'applique au prix.}
\end{solution}

\begin{solution}{iq:iv2:rates-and-fx:4}
Le DV01 vaut 8\,000~\$, et vous gagnez environ 200\,000~\$.

$10\,000\,000 \times 8 \times 0.0001 = 8\,000$, et $25 \times 8\,000 = 200\,000$. Exprimez le DV01 en dollars, car la duration modifiée seule ne dit pas combien vous détenez.

\iqwrong{Oublier qu'une baisse des rendements est un gain pour le détenteur de l'obligation.}
\end{solution}

\begin{solution}{iq:iv2:rates-and-fx:5}
Il vaut environ 5.01~\%. Pour le figer, l'emprunteur achète un zéro-coupon à un an et vend un zéro-coupon à deux ans, ce qui le laisse recevoir de la trésorerie dans un an et en payer dans deux ans.

$1.04^2/1.03 = 1.0501$. Acheter le zéro à un an avec le produit de la vente du zéro à deux ans ne coûte rien aujourd'hui. Le flux de l'année 1 est une entrée, et celui de l'année 2 est ce montant capitalisé au taux de 5.01~\%.

\iqwrong{Faire la moyenne pour obtenir 5~\%. La forme géométrique donne 5.01~\%, et sur des écarts plus longs la différence s'accroît.}
\end{solution}

\begin{solution}{iq:iv2:rates-and-fx:6}
Vendez à découvert 5.625 millions de dollars du dix ans. Vous restez exposé à la forme de la courbe~: le cinq ans et le dix ans ne bougent pas forcément du même montant.

Le DV01 de votre position longue vaut $10\,000\,000 \times 4.5 \times 0.0001 = \$4\,500$. La position courte doit porter ce même DV01~: $4\,500/(8 \times 0.0001) = 5\,625\,000$.

\iqwrong{Couvrir le même notionnel. Vous resteriez court de 3\,500~\$ de DV01, soit un pari important sur les taux.}
\end{solution}

\begin{solution}{iq:iv2:rates-and-fx:7}
Vous gagnez si les taux montent. C'est comme être vendeur d'une obligation à cinq ans à 4~\%, financée par un titre à taux variable. Si le taux de swap monte de 10 points de base, vous gagnez environ 45\,000~\$.

La jambe variable vaut à peu près le pair, et la jambe fixe est une obligation. Payer le taux fixe revient à être vendeur de cette obligation. D'où $10 \times 4\,500$.

\iqwrong{Croire que la jambe variable a un DV01 élevé. Elle se refixe, donc sa duration ne va que jusqu'à la prochaine fixation.}
\end{solution}

\begin{solution}{iq:iv2:rates-and-fx:8}
Le prix vaut 99~\$. Le rendement est d'environ 4.04~\%. Il est supérieur à 4~\% parce que la décote porte sur la valeur nominale, alors que vous touchez votre rendement sur le prix, plus bas.

$100(1 - 0.04 \times 90/360) = 99$. Le rendement vaut $(1/99)(360/90) = 4.04~\%$.

\iqwrong{Traiter un taux d'escompte comme un rendement. Sur un bon du Trésor ils sont proches, mais sur un titre long l'écart grandit.}
\end{solution}

\begin{solution}{iq:iv2:rates-and-fx:9}
Le terme vaut environ 1.1214, soit environ 214 points au-dessus du comptant. Non, il ne le prédit pas. Il est fixé par l'écart de taux de façon qu'aucun arbitrage ne soit possible.

$1.10 \times 1.05/1.03 = 1.12136$. La devise à taux plus bas se traite à une prime à terme. Sinon, emprunter l'une, prêter l'autre, en couverture, rapporterait de l'argent sans risque.

\iqwrong{Lire un terme supérieur au comptant comme une prévision de hausse. Une longue série de tests montre que les taux à terme sont de mauvais prévisionnistes.}
\end{solution}

\begin{solution}{iq:iv2:rates-and-fx:10}
L'EURJPY devrait valoir 165. À 166, vendre de l'EURJPY à ce cours, acheter de l'EURUSD et acheter de l'USDJPY, pour 1 yen par euro. Vérifier d'abord que les trois cotations sont actives, ce que coûtent les fourchettes acheteur-vendeur et la date de règlement de chacune.

$1.10 \times 150 = 165$. Vendre un euro pour 166 yens, le racheter avec 1.10 dollar, puis acheter ces dollars pour 165 yens fait gagner 1 yen.

\iqwrong{Utiliser le prix milieu sur les trois jambes. Franchir trois fourchettes absorbe généralement l'écart qui semble exister à l'écran.}
\end{solution}

\begin{solution}{iq:iv2:rates-and-fx:11}
C'est une option de vente en dollars sur 1.1 million de dollars, à un prix d'exercice d'environ 0.909 euro par dollar.

Un call sur 1 million d'euros à 1.10 vous donne le droit de payer 1.1 million de dollars et de recevoir 1 million d'euros. C'est le droit de vendre 1.1 million de dollars à $1/1.10$ euro par dollar.

\iqwrong{Coter un notionnel d'un million de dollars. Le notionnel de l'autre côté est le premier multiplié par le prix d'exercice.}
\end{solution}

\begin{solution}{iq:iv2:rates-and-fx:12}
Il doit valoir environ 1.1214 pour que vous soyez à l'équilibre, comme le terme. Le risque d'un carry trade est que la devise bouge. Dans un an, l'euro doit monter d'environ 1.9~\% pour effacer votre gain de 2~\%.

Empruntez 1 euro, vendu 1.10~\$. Prêtez à 5~\%, pour obtenir 1.155. Vous devez 1.03 euro. D'où $1.155/1.03 = 1.1214$.

\iqwrong{Parler d'argent gratuit, puisque le terme ne prédit pas en moyenne le mouvement. Le risque a une asymétrie~: les pertes importantes surviennent lors des débouclages.}
\end{solution}

\begin{solution}{iq:iv2:rates-and-fx:13}
Une hausse coûte environ 7.6~\%, et une baisse rapporte environ 8.4~\%. Le gain est plus important. On apprécie cela, car on gagne davantage qu'on ne perd pour le même mouvement.

$-8 \times 1~\% + \tfrac12 \times 80 \times (1~\%)^2 = -7.6~\%$. Le terme de convexité a le même signe dans les deux sens.

\iqwrong{Croire que c'est gratuit. Une courbe valorise la convexité à travers des rendements plus bas sur les obligations longues. La convexité est une position qui a un coût.}
\end{solution}

\begin{solution}{iq:iv2:rates-and-fx:14}
Les facteurs d'actualisation valent $D_1 = 0.97087$ et $D_2 \approx 0.92420$. Le taux zéro-coupon à deux ans vaut environ 4.02~\%, et le terme environ 5.05~\%.

Valorisation au pair à deux ans~: $0.04D_1 + 1.04D_2 = 1$. Donc $D_2 = (1 - 0.04 \times 0.97087)/1.04$. Le taux zéro-coupon vaut $D_2^{-1/2} - 1$, et le terme vaut $D_1/D_2 - 1$.

Ainsi, sur une courbe croissante, le taux zéro-coupon est au-dessus du taux au pair, et le terme est au-dessus des deux.

\iqwrong{Prendre le taux au pair à deux ans pour le taux zéro-coupon à deux ans. L'écart, qui vaut ici 2 points de base, croît avec l'échéance et avec la pente.}
\end{solution}

\begin{solution}{iq:iv2:rates-and-fx:15}
Le taux au pair à trois ans vaut environ 4.93~\%. Le swap à deux ans débutant dans un an vaut environ 5.99~\%. La jambe variable ne demande aucun modèle, puisqu'elle vaut $D_{\text{début}} - D_{\text{fin}}$, par réplication.

Les facteurs d'actualisation valent 0.97087, 0.92456 et 0.86384. Le taux au pair vaut $(1 - D_3)/(D_1 + D_2 + D_3)$. Le taux de swap à terme vaut $(D_1 - D_3)/(D_2 + D_3)$.

Emprunter à la fin et prêter au début de chaque période, et l'on obtient chaque paiement variable sans risque. Cela vaut pour tout modèle de taux, dans le cas simple d'une seule courbe.

\iqwrong{Prévoir les paiements variables puis les actualiser. Si l'on le fait avec les termes, on obtient la même réponse, mais cela masque pourquoi aucun modèle n'est utilisé.}
\end{solution}

\begin{solution}{iq:iv2:rates-and-fx:16}
Le barbell détient environ 41~\% sur le dix ans et 59~\% sur le deux ans. Sa convexité vaut 36.5, contre 25.0 pour le bullet. Taux en baisse de 100~: le barbell gagne 4.64 et le bullet 4.58. Taux en hausse de 100~: le barbell perd 4.28 et le bullet 4.33. Sur une courbe réelle, le bullet offre souvent un rendement plus élevé, ce qui compense son manque de convexité.

Les pondérations viennent de la duration~: 1.886, 4.452 et 8.111. Donc $w_{10} = (4.452 - 1.886)/(8.111 - 1.886) = 0.412$.

Sur une courbe plate, le barbell gagne quoi qu'il arrive, ce qui serait un arbitrage. Une courbe plate ne peut donc pas rester plate dans un monde avec volatilité, et une courbe réelle fait payer la convexité en accordant au cinq ans un rendement plus élevé. C'est pourquoi le trade a un coût de carry.

\iqwrong{Appeler le barbell une convexité gratuite. Il faut aussi comparer le carry et son exposition à un changement de pente.}
\end{solution}

\begin{solution}{iq:iv2:rates-and-fx:17}
On achète environ 41.2~\$ du dix ans et 58.8~\$ du deux ans. On gagne environ 44.5 cents pour 100~\$ si le rendement à cinq ans monte de 10 points de base. On parie que le cinq ans devient meilleur marché que les ailes~: il s'agit d'un butterfly\index{butterfly (rates)}.

Avec $a + b = 100$ et $1.886a + 8.111b = 445.2$, on obtient $b = 41.2$. Le gain vaut $4.452$ cents par point de base sur la position courte, pour 10 points de base.

\iqwrong{Dire que le trade est sans risque, parce qu'il n'a pas de DV01. Il garde une exposition à la pente si les ailes ne bougent pas comme prévu, et il porte une position de convexité.}
\end{solution}

\begin{solution}{iq:iv2:rates-and-fx:18}
Vendre à découvert 10~millions de dollars de dix ans, dont le DV01 vaut 8\,000~\$, et acheter environ 42.1~millions de dollars de deux ans. Si le deux ans baisse de 5 et le dix ans monte de 5, on gagne 80\,000~\$. Si les deux montent de 10, on ne gagne rien.

$8\,000/190 = 42.1$. Chaque jambe rapporte $8\,000 \times 5$.

Un steepener gagne quand la pente augmente, et il n'a pas de duration. Il porte néanmoins du carry~: être acheteur du deux ans et vendeur du dix ans coûte généralement du carry si la courbe monte.

\iqwrong{Utiliser des notionnels égaux. C'est une très forte position vendeuse en duration, et elle a peu à voir avec la pente.}
\end{solution}

\begin{solution}{iq:iv2:rates-and-fx:19}
La régression donne environ 5\,250~\$ de DV01 sur le deux ans, et non 8\,000~\$. La volatilité résiduelle vaut environ 53.6 points de base par an, contre environ 60.2 avec la couverture neutre en DV01.

En régressant la variation du dix ans sur celle du deux ans, on obtient $\beta = 0.7 \times 75/80 = 0.656$. Donc $8\,000 \times 0.656 = 5\,250$. Le résidu vaut $75\sqrt{1 - 0.49} = 53.6$. Le spread neutre en DV01 présente
\[
\sqrt{80^2 + 75^2 - 2 \times 0.7 \times 80 \times 75} = 60.2.
\]

Cela représente environ 27.6~millions de dollars de deux ans, à 190~\$ par million, contre 42.1~millions de dollars. Mais le bêta se déplace avec le régime, et les régressions sur des niveaux de volatilité peuvent induire en erreur. Il faut donc la surveiller et la recalibrer.

\iqwrong{Croire que le DV01 est la seule couverture correcte. Il ne neutralise qu'un déplacement parallèle, et les déplacements réels ne sont pas parallèles.}
\end{solution}

\begin{solution}{iq:iv2:rates-and-fx:20}
Elle vaut environ 1.91~\$ pour 100~\$, par les deux méthodes.

Comme obligation, $4/1.03 + 104/1.03^2 = 101.91$, diminué de 100 pour la jambe variable. Comme écart~: $(4 - 3) \times (1/1.03 + 1/1.03^2) = 1.913$.

\iqwrong{Prendre l'écart de 1~\% multiplié par deux ans, soit 2~\$. Les deux paiements sont actualisés.}
\end{solution}

\begin{solution}{iq:iv2:rates-and-fx:21}
Il vaut environ 13 points de base, et le taux du contrat à terme est au-dessus de celui du terme.

La règle donne $\tfrac12\sigma^2T_1T_2 = 0.5 \times 0.0001 \times 5 \times 5.25 = 13.1$ points de base. Un contrat à terme est réévalué chaque jour. Son détenteur gagne quand les taux baissent et réinvestit à un taux bas~; il perd quand les taux montent et se finance à un taux élevé. Comme cela désavantage la position longue, le prix du contrat à terme est plus bas, ce qui signifie que son taux est plus élevé.

\iqwrong{Ignorer cet effet à l'échéance lointaine. Il croît avec le carré de la maturité, et au bout de cinq ans il dépasse la fourchette.}
\end{solution}

\begin{solution}{iq:iv2:rates-and-fx:22}
La troisième obligation est la moins chère à livrer. Sa base brute est la plus faible, à 0.6, contre 1.0 et 0.8.

Le coût de livraison correspond au prix de l'obligation diminué du prix de facturation~: $93 - 92 = 1.0$, $110.8 - 110 = 0.8$ et $125.6 - 125 = 0.6$.

Le facteur de concordance est le prix à un rendement donné. Loin de ce rendement, les prix des obligations ne varient pas proportionnellement. Lorsque les rendements sont élevés, les obligations à longue duration baissent davantage et deviennent les moins chères à livrer. Lorsqu'ils sont bas, ce sont les obligations à courte duration qui le deviennent.

\iqwrong{Choisir l'obligation au prix le plus bas. Le prix d'une obligation doit être comparé au prix de facturation.}
\end{solution}

\begin{solution}{iq:iv2:rates-and-fx:23}
Le carry vaut environ 1~\%, le roll-down environ 0.72~\%, soit environ 1.72~\% au total. Le rendement à quatre ans peut monter d'environ 48 points de base, depuis 3.8~\%, jusqu'à environ 4.28~\%, avant de perdre de l'argent.

Le carry est le rendement diminué du financement~: $4 - 3$. Au fil de l'année, l'obligation devient une obligation à quatre ans évaluée à 3.8~\%, au lieu de 4~\%, si bien que son prix gagne $0.2 \times 3.6 = 0.72$. Le point mort vaut $1.72/3.6 = 0.48$.

\iqwrong{Oublier le roll. Sur une courbe pentue, il peut dépasser le carry. Et il n'apparaît que si la courbe reste inchangée au cours du temps.}
\end{solution}

\begin{solution}{iq:iv2:rates-and-fx:24}
La moyenne vaut environ 4.26~\%, et l'écart-type environ 0.93~\%. La probabilité d'un taux négatif est d'environ 0.0002~\%. Non, cela ne le rend pas inutile. Ce point ne représente souvent qu'une petite part des prix du modèle.

La moyenne vaut $5 + (3 - 5)e^{-1}$. La variance vaut $\sigma^2(1 - e^{-2\kappa T})/(2\kappa)$. La probabilité vaut $\Phi(-4.26/0.93)$.

Un modèle normal a ses mérites~: des formules fermées, et une volatilité qui correspond à la manière dont les taux ont tendance à évoluer en points de base. Des taux négatifs se sont produits sur certains marchés, si bien que cet aspect n'est pas toujours un défaut.

\iqwrong{Rejeter un modèle à cause d'un trait du modèle sans effet sur le prix considéré. Il faut d'abord tester cet effet.}
\end{solution}

\begin{solution}{iq:iv2:rates-and-fx:25}
Cela représente environ 16.7 points de base. Les rendements longs baissent, puisqu'un taux aléatoire rend le facteur d'actualisation convexe par rapport à ce taux.

Le taux intégré a une variance de $\sigma^2T^3/3$. Donc $E[e^{-\int r}] = e^{-f T + \sigma^2T^3/6}$. Le rendement vaut $f - \sigma^2T^2/6$. Pour $T = 10$, $\sigma^2 = 0.0001$~: $0.0001 \times 100/6 = 0.00167$. La courbe à terme de ce modèle vaut $f - \sigma^2T^2/2$~; une courbe à terme parfaitement plate donnerait un rendement zéro-coupon égal exactement à son niveau.

C'est à l'origine de la valeur de convexité que portent les obligations longues. Ainsi, les rendements ont tendance à passer sous les taux anticipés aux longues échéances.

\iqwrong{Croire que la volatilité n'affecte que les options. L'obligation est une moyenne de fonctions convexes du taux, donc la volatilité l'affecte aussi.}
\end{solution}

\begin{solution}{iq:iv2:rates-and-fx:26}
Le caplet vaut environ 7.57 points de base du notionnel. La même volatilité vaut environ 80 points de base, et Bachelier donne environ 7.58. À un taux de 2~\%, il vaudrait environ 40 points de base.

La formule de Black sur le taux forward, avec une fraction de période de 0.25 et un facteur d'actualisation de 0.95. À la monnaie, elle vaut $0.95 \times 0.25 \times 0.04 \times (2\Phi(0.1) - 1) = 0.000757$. La forme normale vaut $0.95 \times 0.25 \times 0.008 \times 0.3989$. Pour convertir, on utilise $\sigma_N \approx F\sigma_L$.

\iqwrong{Comparer la volatilité lognormale d'un niveau de taux à l'autre. Un passage de 20~\% à 40~\% alors que le taux est divisé par deux peut correspondre à la même volatilité normale. Les traders cotent en volatilité normale.}
\end{solution}

\begin{solution}{iq:iv2:rates-and-fx:27}
Elle vaut environ 1.44~\% du notionnel. L'annuité la multiplie, car chaque point de base de taux de swap est payé à toutes les dates futures. Un point de base de volatilité vaut environ 1.8 point de base du notionnel.

La règle est $A\sigma_N\sqrt{T/(2\pi)} = 4.5 \times 0.008 \times 0.3989$. Le payoff vaut $A(S - K)^+$ à l'échéance. En prenant l'annuité comme numéraire, le taux de swap est une martingale. Pour la sensibilité, $4.5 \times 0.0001 \times 0.3989 = 0.018~\%$.

\iqwrong{Utiliser un facteur d'actualisation à la place de l'annuité. Cela revient à traiter le swap comme ayant un seul paiement.}
\end{solution}

\begin{solution}{iq:iv2:rates-and-fx:28}
Elle vaut $-4.5~\%$ du notionnel. C'est un swap à départ différé à 5~\%, et aucun modèle n'est nécessaire.

À l'échéance, les deux options paient $A(S - K)^+ - A(K - S)^+ = A(S - K)$. Donc aujourd'hui, en unités d'annuité, elle vaut $A(F - K) = 4.5 \times (4 - 5)$.

\iqwrong{Coter une volatilité pour ce produit. La volatilité n'affecte pas la différence.}
\end{solution}

\begin{solution}{iq:iv2:rates-and-fx:29}
Le taux implicite ressort à environ 5.29~\%, soit environ 29 points de base au-dessus du taux cash. Prêter des dollars via le marché des swaps rapporte donc davantage qu'en cash. C'est une base, qui ne se referme pas, car l'arbitrage exige du bilan et du crédit, dont le coût est bien réel.

Le taux implicite vaut $1.1245/1.10 \times 1.03 - 1 = 5.294~\%$.

Le trade consiste donc à emprunter des dollars cash à 5~\%, à les vendre contre des euros, à placer ces euros à 3~\% et à acheter des dollars à terme. Cela verrouille environ 29 points de base. On a longtemps tenu la parité des taux d'intérêt couverte pour presque une loi~; depuis la crise de 2008, elle est systématiquement mise en défaut. Pour une banque, la contrainte tient à ses règles de fonds propres et au coût de son propre bilan.

\iqwrong{Traiter cela comme un gain sans risque. Qui peut le faire, à quel coût de financement et avec quel bilan, voilà les vraies questions.}
\end{solution}

\begin{solution}{iq:iv2:rates-and-fx:30}
Le forward vaut environ 1.1222. Le call vaut environ 0.0536~\$, et le put environ 0.0324~\$. Le delta spot vaut 0.581, et le delta forward 0.599.

La forme de Garman-Kohlhagen\index{Garman-Kohlhagen} traite le taux étranger comme un rendement de dividende. $F = Se^{(r_d - r_f)T}$. Donc $C = e^{-r_d}[F\Phi(d_1) - K\Phi(d_2)]$, avec $d_1 = (\ln(F/K) + 0.005)/0.1$.

Le delta spot vaut $e^{-r_f}\Phi(d_1)$. C'est la quantité d'euros que l'on détient aujourd'hui. Le delta forward, $\Phi(d_1)$, correspond à ce que l'on détient en forwards.

Parité~: $C - P = e^{-r_d}(F - K) = 0.0211$.

\iqwrong{Utiliser la formule d'une action avec un seul taux. Une devise étrangère rapporte un intérêt, donc son forward ne se situe pas au spot capitalisé au seul taux en dollars.}
\end{solution}

\begin{solution}{iq:iv2:rates-and-fx:31}
Elle vaut environ 4.87 centimes d'euro par euro. Cela fait 4.87~\% du notionnel en euros et 4.87~\% du notionnel en dollars également, puisque le strike et le spot sont ici identiques. Valorisée côté dollar, elle doit utiliser la même volatilité de 10~\%, car le logarithme d'un taux et celui de son inverse ont la même variance.

$0.0536/1.10 = 0.0487$. En pourcentage du notionnel en dollars, cela vaut $0.0536/K$. Avec $K = S$, les deux grandeurs sont égales. Le put en dollars sur $X = 1/S$, de strike $1/K$, par dollar, est un put européen de forward $F' = 1/F$. Les taux échangent leurs rôles. Cela fait environ 0.0443 euro par dollar, soit 0.0487 euro pour 1.10~\$, ou 0.0536~\$ en valeur.

\iqwrong{Croire qu'un put en USD sur EURUSD a un smile différent. C'est la même option. C'est la convention, et non le risque, qui diffère.}
\end{solution}

\begin{solution}{iq:iv2:rates-and-fx:32}
La volatilité du call vaut 9.5~\%, et celle du put 11.5~\%. Les strikes sont environ 1.199 et 1.048.

La volatilité du call vaut à la monnaie + butterfly + la moitié du risk reversal, et celle du put vaut à la monnaie + butterfly $-$ la moitié du risk reversal. Pour trouver le strike, on résout $e^{-r_f}\Phi(d_1) = 0.25$ pour le call. Donc $\Phi(d_1) = 0.2576$, $d_1 = -0.651$, et $K = Fe^{-d_1\sigma + \sigma^2/2} = 1.1222e^{0.0618 + 0.0045}$.

\iqwrong{Lire les strikes comme fixes. En change, ils ne sont pas cotés directement, et un delta est un strike mobile qui dépend de la volatilité elle-même.}
\end{solution}

\begin{solution}{iq:iv2:rates-and-fx:33}
Elle vaut environ 16.5~\% avec une corrélation de 0.5, et environ 9.5~\% avec une corrélation de $-0.5$.

L'EURJPY est le produit des deux autres, donc $\sigma^2 = \sigma_1^2 + \sigma_2^2 + 2\rho\sigma_1\sigma_2$. Soit $0.01 + 0.0081 + 0.009 = 0.0271$, dont la racine carrée vaut 16.46~\%.

À l'inverse, trois volatilités implicites d'options donnent une corrélation implicite. Les traders comparent celle-ci à la corrélation réalisée. Une paire sans options liquides peut être valorisée à partir des deux autres, et l'écart de prix est un trade de corrélation.

\iqwrong{Utiliser un signe moins. C'est un produit, pas une différence. Pour un rapport comme l'EURGBP, à partir de l'EURUSD et du GBPUSD, le signe est inversé.}
\end{solution}

\begin{solution}{iq:iv2:rates-and-fx:34}
Le taux composé vaut environ 4.144~\%, et la moyenne simple environ 4.143~\%. Il n'est connu qu'à la fin, ce qui signifie qu'un caplet sur ce taux continue de bouger pendant la période de calcul, si bien que sa volatilité ne s'arrête pas au début.

Chaque jour multiplie par $1 + r_id_i/360$. Le produit vaut 1.0008058, et le taux vaut ce dernier moins 1, multiplié par 360/7.

\begin{dated}{2026-11}{Taux overnight}\label{dat:iv2:rates-and-fx:ois}
En dollars américains, les swaps overnight sont indexés sur le SOFR, avec des paiements annuels, actual/360 sur les deux jambes, et un taux variable composé a posteriori. Ces conventions évoluent~: il faut les vérifier avant de passer un trade. Voir le registre du chapitre.
\end{dated}

Un caplet composé a posteriori a donc une variance totale qui inclut la période de calcul, avec une variance qui décroît au fil de celle-ci. Une règle empirique donne un tiers de la période de calcul comme temps de variance supplémentaire.

\iqwrong{Croire qu'un taux composé est assez proche d'une moyenne pour tous les usages. Sur des périodes plus longues et avec des taux plus élevés, ce n'est plus le cas.}
\end{solution}

\begin{solution}{iq:iv2:rates-and-fx:35}
Le zéro-coupon à cinq ans vaut environ 0.8083, soit un rendement d'environ 4.26~\% en capitalisation continue. Une simulation donne 0.8082. À très longue échéance, le rendement tend vers $\theta - \sigma^2/(2\kappa^2)$, soit un niveau de 4.98~\%.

On écrit $P = A(T)e^{-B(T)r}$. En substituant, on obtient $B' + \kappa B = 1$, donc $B = (1 - e^{-\kappa T})/\kappa$. Et $\ln A = (\theta - \sigma^2/(2\kappa^2))(B - T) - \sigma^2B^2/(4\kappa)$.

Pour la vérification, on simule le taux, on l'intègre par la règle des trapèzes, et l'on moyenne le facteur d'actualisation.

Ainsi, le rendement long est inférieur à $\theta$, de la même valeur de convexité que dans Ho-Lee, mais limité par le retour à la moyenne. C'est la vitesse de retour qui le borne~: sans retour à la moyenne, l'écart croît sans limite.

\iqwrong{Supposer que le rendement tend vers la moyenne du taux. Le terme de convexité le fait baisser, un peu ici.}
\end{solution}

\begin{solution}{iq:iv2:rates-and-fx:36}
La dérive est $\theta(t) = \partial_tf(0, t) + \kappa f(0, t) + \sigma^2(1 - e^{-2\kappa t})/(2\kappa)$. Une courbe plate donne $\kappa \cdot 4~\% + \sigma^2(1 - e^{-2\kappa t})/(2\kappa)$. La constante de Vasicek ne peut pas reproduire une courbe entière, car elle a trop peu de paramètres. Dans une simulation, le modèle restitue le zéro-coupon à dix ans de $e^{-0.4}$.

La dérive est ce qu'exige une courbe pour être ajustée exactement. Le forward dans le modèle est $f(t, T) = f(0, T) + \tfrac12\sigma^2\bigl(B(0, T)^2 - B(t, T)^2\bigr) + e^{-\kappa(T - t)}x_t$, avec $B(t, T) = (1 - e^{-\kappa(T - t)})/\kappa$, $\mathrm{d}x = -\kappa x\,\mathrm{d}t + \sigma\,\mathrm{d}W$ et $x_0 = 0$. Il se déplace donc avec un choc qui s'atténue avec l'échéance, à la vitesse $\kappa$.

On simule $r$ sur de petits pas de temps avec cette dérive, et on l'intègre le long de chaque trajectoire. Avec le terme ci-dessus, le facteur d'actualisation moyen s'accorde au marché à la précision du bruit de simulation près. Sans lui, le modèle manque la courbe.

\iqwrong{Calibrer le retour à la moyenne sur l'historique, et affirmer que cela ajuste la courbe. Dans Hull-White, la vitesse vient des options et la dérive vient de la courbe.}
\end{solution}

\begin{solution}{iq:iv2:rates-and-fx:37}
Le cap vaut environ 59.9 points de base. La swaption vaut environ 60, 52 et 43 pour des corrélations de 1, 0.5 et 0. Elle ne vaut jamais davantage, car une option sur une moyenne ne vaut jamais plus que la moyenne des options. C'est ce qui relie les marchés du cap et des swaptions~: leur écart reflète la corrélation.

Le cap vaut $1.88 \times$ le caplet, soit $0.04(2\Phi(0.1) - 1)$. La swaption vaut $1.88E[(\tfrac12(F_1 + F_2) - K)^+]$. Avec $\rho = 1$, les deux sont égales.

Par convexité, le maximum d'une somme ne dépasse pas la somme des maxima. Une corrélation inférieure à 1 abaisse donc la variance de la moyenne, et cela fait baisser la swaption.

\iqwrong{Calibrer un modèle à un facteur sur les caps et sur les swaptions et s'attendre à un ajustement. Un seul facteur signifie que les taux forward sont presque parfaitement corrélés, de sorte qu'il surévalue les swaptions par rapport aux caps.}
\end{solution}

\begin{solution}{iq:iv2:rates-and-fx:38}
Les obligations bullet valent 104.45, 100.00 et 95.79. Les calls valent 3.86, 1.39 et 0.29. Les obligations callables valent 100.60, 98.61 et 95.50. Leurs durations valent 1.44, 2.60 et 3.76, contre environ 4.3 pour les bullets. Le prix de la callable monte à peine quand les taux baissent, ce qui est une convexité négative\index{convexité négative}. Son détenteur est vendeur de cette option.

À faible taux, l'émetteur rachète, de sorte que l'investisseur récupère le pair, et l'obligation se comporte comme une obligation à un an. À taux élevé, le call ne vaut rien, et l'obligation se comporte comme une obligation bullet. La duration se raccourcit donc quand les taux baissent, ce qui est l'inverse de ce que souhaite le détenteur.

Une obligation hypothécaire a la même forme, car les emprunteurs se refinancent. Un détenteur qui se couvre en vendant davantage d'obligations quand les taux montent alimente la baisse.

\iqwrong{Comparer les rendements sans tenir compte de la valeur de l'option. Le rendement supplémentaire d'une callable est la prime de la vente de ce call.}
\end{solution}

\begin{solution}{iq:iv2:rates-and-fx:39}
Le prix juste est d'environ 97.6, soit environ 2.5 de moins que la plus basse des deux moyennes. Si les taux montent, l'obligation longue est la moins chère à livrer. S'ils baissent, c'est l'obligation courte.

\begin{sloppypar}Les deux obligations valent 100 à 6~\%. À la livraison, le prix de la moins chère est $\min(P_5(y), P_{20}(y))$. Le vendeur détient une option, et elle coûte à l'acheteur.\end{sloppypar}

L'obligation longue a beaucoup plus de duration. Le delta du contrat à terme change donc avec la moins chère à livrer, et il se comporte comme si l'acheteur avait vendu un straddle au vendeur. Le delta usuel, qui ne porte que sur une obligation, sous-estime cet effet.

\iqwrong{Évaluer le contrat à terme à partir de la moins chère à livrer du jour et de son portage seulement. Cela ignore une option réelle, de valeur réelle, lorsque les obligations sont proches l'une de l'autre.}
\end{solution}

\begin{solution}{iq:iv2:rates-and-fx:40}
Les trois composantes expliquent environ 85.3~\%, 10.7~\% et 2.6~\%. La première fait bouger toutes les échéances dans le même sens, la deuxième fait bouger les courtes et les longues en sens opposés, et la troisième fait bouger le milieu à contre-courant des extrémités. Il s'agit du niveau, de la pente et de la courbure\index{composantes principales}.

Prendre la matrice de covariance et calculer ses vecteurs propres. Le premier a pour poids environ 0.50, 0.59, 0.50 et 0.40. Le deuxième a pour poids $-0.74$, $-0.09$, 0.40, 0.53. Le troisième a pour poids $-0.42$, 0.62, 0.19, $-0.63$.

On peut donc couvrir les deux premières en détenant deux instruments qui les annulent, et ne rester exposé qu'à la troisième. C'est un butterfly pondéré par ces charges, à la place d'un butterfly neutre en DV01. Toutefois, les composantes sont estimées sur un échantillon, et leurs poids se déplacent. Il faut donc les réestimer et les surveiller.

\iqwrong{Affirmer que la première composante est un pur décalage parallèle. Ses poids sont proches mais différents, ce qui explique que les couvertures en DV01 fuient.}
\end{solution}

\begin{solution}{iq:iv2:rates-and-fx:41}
L'ajustement exact est d'environ 3.41 points de base, et la règle du premier ordre donne environ 3.08. Il apparaît parce que la date de paiement ne coïncide pas avec la date de paiement naturelle.

Sous la probabilité du paiement en $T + \tau$, $F$ est une martingale. Pour payer en $T$, on change de probabilité avec $P(T, T + \tau)^{-1} = 1 + \tau F_T$. Donc
\[
E^T[F_T] = \frac{E^{T+\tau}[F_T(1 + \tau F_T)]}{1 + \tau F_0} = F_0 + \frac{\tau F_0^2(e^{\sigma^2T} - 1)}{1 + \tau F_0}.
\]
Avec $\tau = 1$, $F_0 = 0.04$ et $\sigma^2T = 0.2$, cela donne 3.41 points de base. La règle, $\tau F^2\sigma^2T/(1 + \tau F)$, donne 3.08.

Taux et actualisation sont liés. Un taux élevé est payé tôt, à un moment où il vaut davantage, et c'est favorable au détenteur.

\iqwrong{Utiliser le forward sans ajustement. Il est faible ici, mais pour un taux à long terme payé tôt, il dépasse la fourchette.}
\end{solution}

\begin{solution}{iq:iv2:rates-and-fx:42}
La règle du premier ordre donne environ 17.6 points de base, et la simulation environ 18.2. L'ajustement relève le taux. Il dépend de tout le smile, car une simulation du payoff est une fonction du taux pour tous les strikes.

Sous la probabilité annuité, le taux de swap $S$ est une martingale. Pour payer à la date de fixing, on change de probabilité~: $E^T[S] = E^A[S/A(S)]\,/\,E^A[1/A(S)]$, où le dénominateur donne $E^T[1] = 1$. Avec une forme à taux plat, $A(y) = (1 - (1 + y)^{-10})/y$. À 4~\%, $A = 8.11$ et $A' = -40.4$. La règle vaut $F^2(e^{\sigma^2T} - 1)(-A'/A) = 0.0016 \times 0.2214 \times 4.98$.

Comme $A$ baisse quand $S$ monte, $1/A$ est maximal quand le taux de swap est élevé. Le payoff présente donc une inclinaison convexe. Sur le marché, l'intégrale porte sur les swaptions à tous les strikes, et un smile qui relève les ailes relève aussi cet ajustement.

\iqwrong{Utiliser le taux de swap forward pour le fixing espéré. Sur des CMS à long terme, l'erreur peut atteindre plusieurs points de base.}
\end{solution}

\begin{solution}{iq:iv2:rates-and-fx:43}
Le quanto\index{quanto} forward vaut environ 100.60 à un an, au-dessus du spot, et une simulation confirme ce résultat. Avec une corrélation positive, il serait inférieur.

L'indice en dollars, $SX$, est un actif négociable, où $X$ est la valeur en dollars d'un yen. Donc, avec des taux nuls, $SX$ n'a pas de dérive sous la probabilité du dollar, et il en va de même pour $X$. Par Itô, $d(SX)/(SX) = (\mu + \rho\sigma_S\sigma_X)\,dt + \ldots$, donc la dérive de l'indice doit être $\mu = -\rho\sigma_S\sigma_X$. On obtient alors $100e^{0.3 \times 0.2 \times 0.1}$.

Pour la vérification, passer par la probabilité du yen. Là, $S$ et $Y = 1/X$ sont des martingales, avec une corrélation de $+0.3$. Le changement vers la probabilité du dollar pondère par $Y$, de sorte que $E^{\text{usd}}[S] = E^{\text{yen}}[SY]/E^{\text{yen}}[Y]$. La simulation donne 100.60.

Le signe a un sens. Un investisseur en dollars qui détient l'indice réel est pénalisé lorsque le yen baisse alors que l'indice monte, ce qu'indique la corrélation ici. Le quanto est exempt de ce risque, et il vaut davantage.

\iqwrong{Supposer le signe sans préciser ce que mesure $X$. Le signe de la corrélation s'inverse quand on cote le taux dans l'autre sens, donc il faut d'abord le définir.}
\end{solution}

\begin{solution}{iq:iv2:rates-and-fx:44}
Sans prime dans le delta, le strike vaut $Fe^{\sigma^2T/2} \approx 1.1278$. Avec la prime incluse, il vaut $Fe^{-\sigma^2T/2} \approx 1.1166$. Les paires dont la prime est payée en devise étrangère utilisent plutôt la seconde, car la prime est elle-même une position dans cette devise.

Pour que le delta d'un call soit égal à $-1$ fois celui d'un put, on pose $\Phi(d_1) = \tfrac12$, donc $d_1 = 0$. Avec la forme ajustée de la prime, le delta vaut $(K/F)\Phi(d_2)$, et les deux jambes s'annulent quand $d_2 = 0$.

Si on les confond, la cotation à 25 delta est fixée sur le mauvais strike, et le risk reversal qui en découle est faux. Le smile d'un desk s'ajuste mal, et sa couverture est fausse.

\iqwrong{Traiter la parité à la monnaie comme un point unique. Il en existe plusieurs, et chaque marché a sa convention.}
\end{solution}

\begin{solution}{iq:iv2:rates-and-fx:45}
Le forward vaut environ 6.0~\% juste avant deux ans, et environ 3.6~\% juste après. Il vaut 3.2~\% juste avant trois ans, et 4.4~\% juste après. Avec des facteurs d'actualisation log-linéaires, il est plat sur chaque période à 5.0~\%, 3.4~\% et 4.8~\%. On se couvrirait avec le second, car les dents de scie du premier n'ont aucun sens.

Pour un taux zéro linéaire, $f(t) = z(t) + tz'(t)$. Entre un et deux ans, la pente vaut 1~\%, donc à deux ans, $4~\% + 2 \times 1~\% = 6~\%$. Juste après, la pente vaut $-0.2~\%$, donc $4 - 0.4 = 3.6~\%$.

Avec des forwards constants par morceaux, le taux est la différence log des facteurs d'actualisation. La modification d'une donnée ne touche qu'un segment local. Ses deltas restent donc dans les compartiments voisins, alors qu'une spline cubique peut répartir un choc vers des points éloignés.

\iqwrong{Choisir un schéma d'après l'allure d'une courbe de rendement. Le test porte sur ce qui arrive à ses forwards et à son risque.}
\end{solution}

\begin{solution}{iq:iv2:rates-and-fx:46}
La volatilité des prix de l'obligation vaut environ 3.56~\%. Le call vaut environ 1.42~\%. La swaption receveuse vaut environ 1.44~\%. Les deux valeurs sont presque identiques, car chacune a la même exposition aux mouvements de taux sur un an. La swaption a l'exposition la plus forte à son annuité.

$4.45 \times 0.8~\% = 3.56~\%$. La règle donne $0.3989 \times 3.56~\% = 1.42~\%$, et celle de la swaption vaut $4.5 \times 0.8~\% \times 0.3989 = 1.44~\%$. Un call sur obligation gagne quand les taux baissent, comme un receveur.

Leurs différences~: une obligation à cinq ans glisse vers quatre ans à l'échéance, donc sa duration forward à l'échéance est plus courte. Il faut donc utiliser la duration forward. Le prix de l'obligation forward est aussi financé par la pension livrée. Les deux ont donc un portage et un risque de special différents.

\iqwrong{Utiliser la duration d'aujourd'hui pour une option sur obligation, alors que l'obligation a vieilli à l'échéance. Le vieillissement d'un an retire environ un point de duration.}
\end{solution}

\begin{solution}{iq:iv2:rates-and-fx:47}
Le taux de repo implicite vaut 3~\%. On achète donc l'obligation en la finançant à 2.5~\%, et on vend le contrat à terme, pour fixer environ 0.25~\% sur un semestre.

Le taux vaut $(99.5 + 2 - 100)/100 \times (1/0.5) = 3\%$. C'est ce que rapporte la détention de l'obligation jusqu'à sa livraison. Son financement coûte moins cher. C'est un cash-and-carry\index{basis trade}.

Ce qui peut gâcher l'opération~: un changement de la moins chère à livrer (on détient alors la mauvaise obligation)~; les options de livraison du vendeur, qu'on détient sans toujours en tirer de la valeur~; une variation du taux de repo avant la livraison~; et, en situation de stress, la baisse de l'offre de repo, qui rend l'effet de levier dangereux. Ces opérations rapportent peu par unité, et elles recourent donc à un fort effet de levier.

\iqwrong{Dire que c'est sans risque. La base a déjà explosé par le passé, quand le financement s'est resserré, et les tailles ont vite grossi.}
\end{solution}

\begin{solution}{iq:iv2:rates-and-fx:48}
Le strike vaut environ 4.08~\%, en capitalisation annuelle. La swaption vaut environ 1.338~\% du notionnel. Une simulation donne 1.335, avec une erreur de 0.002.

Une swaption payeuse est un put sur une obligation à coupons de strike 1. Dans un modèle à un facteur, tous les prix des obligations à l'échéance baissent quand $r$ monte. Il existe donc un seul $r^*$ pour lequel l'obligation à coupons vaut 1. Chaque zéro-coupon a alors son propre strike $X_i = P(1, T_i, r^*)$, et l'option sur la somme égale la somme des options sur les zéro-coupons. C'est l'astuce de Jamshidian\index{décomposition de Jamshidian}. Le put de chaque zéro-coupon a une formule fermée.

Dans un modèle à deux facteurs, le fait que la valeur de l'obligation soit inférieure au pair ne correspond pas à un seul niveau. L'ensemble est une courbe dans l'espace des états, et l'ordre des prix des obligations n'est pas fixe, de sorte qu'il ne se scinde pas.

Pour la simulation, on utilise la probabilité forward à un an, dans laquelle la moyenne du taux court à l'échéance est le forward. On actualise ensuite avec $P(0, 1)$.

\iqwrong{Appliquer une formule de Black à volatilité plate à un prix de modèle. La volatilité implicite de swaption d'un modèle est un résultat, et elle varie avec la vitesse.}
\end{solution}

\begin{solution}{iq:iv2:rates-and-fx:49}
Elle réplique le vega, le vanna et le volga de la nouvelle option à l'aide de trois options liquides. Le prix de la nouvelle option est son prix à volatilité plate, plus les pondérations multipliées par le coût du smile de chaque pilier. On obtient 10.41~\% à 1.10 et 9.57~\% à 1.25. Les piliers sont retrouvés exactement, car les pondérations y valent 1 et 0.

Les pondérations résolvent un système trois par trois. Celui-ci admet une forme fermée~:
\[
x_1(K) = \frac{\text{vega}(K)}{\text{vega}(K_1)}\frac{\ln(K_2/K)\ln(K_3/K)}{\ln(K_2/K_1)\ln(K_3/K_1)},
\]
et de même avec les indices permutés. En $K = 1.25$, les pondérations valent 0.33, $-0.96$ et 1.56.

La formule du premier ordre donne de bonnes volatilités au voisinage des piliers. En dehors de la zone des 25~delta, les pondérations croissent, et le smile peut se comporter mal. En 1.30, elle donne 9.97~\%, en hausse rapide par rapport à 9.57~\%. Il faut donc vérifier l'absence d'arbitrage avant de l'utiliser dans les ailes, et recourir à un modèle ou à un plancher dans la queue lointaine.

\iqwrong{Prendre le vanna-volga pour un modèle. C'est une règle d'interpolation assortie d'une histoire de couverture. Elle n'a pas de dynamique, et elle peut se tromper sur les produits exotiques.}
\end{solution}

\begin{solution}{iq:iv2:rates-and-fx:50}
Les européennes valent environ 1.14, 1.18, 0.95 et 0.54 en pourcentage du notionnel, et la bermudéenne environ 1.585. La valeur de l'option d'échange par rapport à la meilleure européenne vaut donc environ 0.40, soit 34~\%. Avec une vitesse de 0.3 et la même volatilité, toutes les valeurs baissent~: la meilleure européenne vaut 0.76, la bermudéenne 1.105, et la valeur de l'option d'échange représente 45~\% de la meilleure européenne. Une vitesse plus élevée augmente cette part, car les taux de swap aux différentes dates deviennent moins corrélés.

Notons $x$ l'état gaussien sans dérive. Sur une grille en $x$, on remonte de l'année 4 à l'année 1. À chaque date, on prend le plus grand de la valeur de la détention et de la valeur d'exercice du swap. Pour passer de l'année $t + 1$ à l'année $t$, on intègre sous la probabilité $(t + 1)$-forward, dont la transition est gaussienne avec un décalage connu, puis on multiplie par le facteur d'actualisation $P(t, t + 1, x)$.

Pour le premier contrôle, la grille reprice un zéro-coupon~: elle redonne $P(0, 3)$ à $10^{-7}$ près. Pour le second, la première européenne est en accord avec la forme fermée de Jamshidian.

Une bermudéenne vaut au moins sa meilleure européenne. Avec une vitesse plus élevée, les chocs s'éteignent plus vite, de sorte que les taux de swap ultérieurs dépendent moins des précédents. Le détenteur dispose alors de davantage de choix réels, et le droit supplémentaire vaut davantage par rapport à une européenne prise isolément. C'est pourquoi les prix des bermudéennes servent à lire la vitesse.

\iqwrong{Calibrer le modèle sur un seul jeu d'européennes, puis coter une bermudéenne. La valeur de l'option d'échange dépend de la vitesse, et les européennes ne la déterminent pas.}
\end{solution}
